\chapter{Perímetro, área e volume}\label{ch:g6:measure}

Que comprimento tem a cerca, que tamanho tem o campo, quanta água cabe
no tanque? Três perguntas diferentes, três grandezas diferentes ---
\omterm{def:g6:measure:perimeter}{perímetro}, \omterm{def:g6:measure:area}{área}, \omterm{def:g6:measure:volume}{volume} ---, cada uma com as suas unidades. Confundi-las
é o erro mais comum da geometria; este capítulo põe ordem nisso de uma
vez por todas.

\section{Comprimentos e perímetro}

\begin{definition}[Perímetro]\label{def:g6:measure:perimeter}
O \emph{perímetro}\index{perímetro} de uma figura é o comprimento total
do contorno dela. Ele é medido em unidades de comprimento: milímetros
(mm), centímetros (cm), metros (m), quilômetros (km), com
\[
1 \text{ km} = 1000 \text{ m}, \qquad
1 \text{ m} = 100 \text{ cm}, \qquad
1 \text{ cm} = 10 \text{ mm}.
\]
\end{definition}

\begin{example}\label{ex:g6:measure:perimeter}
Um retângulo de comprimento $L$ e largura $w$ tem \omterm{def:g6:measure:perimeter}{perímetro}
\[
P = L + w + L + w = 2 \times (L + w).
\]
Para um retângulo de $7$ cm por $4$ cm: $P = 2 \times (7 + 4) =
2 \times 11 = 22$ cm. Um quadrado de lado $c$ tem \omterm{def:g6:measure:perimeter}{perímetro} $4c$.
\end{example}

\begin{proposition}[Comprimento de uma circunferência]\label{prop:g6:measure:circle}
O \omterm{def:g6:measure:perimeter}{perímetro} (ou \emph{comprimento da circunferência}) de um círculo de
\omterm{def:g3:shapes:shapes}{raio} $r$ é
\[
P = 2 \pi r ,
\]
em que $\pi \approx 3.14$ é o mesmo número para todos os círculos.
\end{proposition}
\admitted

\begin{example}
Um lago circular tem \omterm{def:g3:shapes:shapes}{raio} $5$ m. A borda dele mede
$2 \times \pi \times 5 = 10\pi \approx 31.4$ m. Guarde o valor exato
$10\pi$ o máximo possível; arredonde só no fim.
\end{example}

\section{Áreas}

\begin{definition}[Área]\label{def:g6:measure:area}
A \emph{área}\index{área} de uma figura mede a superfície que ela
cobre: quantos quadrados unitários cabem dentro dela. Unidades: cm$^2$
(um quadrado de lado $1$ cm), m$^2$, km$^2$, \dots\ Cuidado:
\[
1 \text{ m}^2 = 100 \times 100 \text{ cm}^2 = 10\,000 \text{ cm}^2
\]
--- um metro quadrado é um quadrado de $100$ cm por $100$ cm, então
cada degrau da escada de unidades vale $100$, e não $10$.
\end{definition}

\begin{omfigure}
\begin{tikzpicture}[scale=0.68]
  \draw[gray!50, thin] (0,0) grid (7,4);
  \draw[very thick, omDef] (0,0) rectangle (7,4);
  \fill[omThm!25] (0,3) rectangle (1,4);
  \draw[omThm, thick] (0,3) rectangle (1,4);
  \node[font=\small, omThm] at (2.6,3.5) {$1$ quadrado unitário};
  \node[font=\small] at (3.5,-0.5) {$7 \times 4 = 28$ quadrados unitários};
\end{tikzpicture}

{\small A \omterm{def:g6:measure:area}{área} de um retângulo: $7$ colunas de $4$ quadrados unitários
cada uma, logo $7 \times 4 = 28$ quadrados ao todo. É por isso que a
\omterm{def:g6:measure:area}{área} é comprimento $\times$ largura.}
\end{omfigure}

\begin{proposition}[Fórmulas básicas de área]\label{prop:g6:measure:formulas}
\begin{center}
\renewcommand{\arraystretch}{1.5}
\begin{tabular}{l|l}
figura & \omterm{def:g6:measure:area}{área} \\
\hline
retângulo ($L \times w$) & $A = L \times w$ \\
quadrado (lado $c$) & $A = c^2$ \\[4pt]
\omterm{def:g6:shapes:triangles}{triângulo retângulo} (catetos $a$, $b$) & $A = \dfrac{a \times b}{2}$ \\[4pt]
disco (\omterm{def:g3:shapes:shapes}{raio} $r$) & $A = \pi r^2$ \\
\end{tabular}
\end{center}
\end{proposition}

\begin{proof}[Demonstração para o triângulo retângulo]
Duas cópias de um \omterm{def:g6:shapes:triangles}{triângulo retângulo} de catetos $a$ e $b$, coladas
pela hipotenusa, formam um retângulo $a \times b$. Logo a \omterm{def:g6:measure:area}{área} do
triângulo é a \omterm{def:g3:division:half}{metade} da do retângulo: $\frac{ab}{2}$. (A fórmula do
retângulo é a contagem de quadrados unitários feita acima; a fórmula do
disco é admitida, veja o \cref{ch:g7:areas}.)
\end{proof}

\begin{omfigure}
\begin{tikzpicture}[scale=0.75]
  \draw[thick, fill=omDef!15] (0,0) -- (4,0) -- (0,2.5) -- cycle;
  \draw[thick, omExo, dashed] (4,0) -- (4,2.5) -- (0,2.5);
  \draw[thin] (0.3,0) -- (0.3,0.3) -- (0,0.3);
  \node[below, font=\small] at (2,-0.1) {$a$};
  \node[left, font=\small] at (0,1.25) {$b$};
  \node[font=\small, omExo] at (2.8,1.7) {segunda cópia};
\end{tikzpicture}

{\small Um \omterm{def:g6:shapes:triangles}{triângulo retângulo} é \omterm{def:g3:division:half}{metade} de um retângulo: daí a fórmula
$\frac{a\times b}{2}$.}
\end{omfigure}

\begin{example}[Figuras compostas]\label{ex:g6:measure:composite}
Um cômodo em forma de L é um retângulo de $6$ m $\times$ $4$ m com um
canto quadrado de $2$ m $\times$ $2$ m retirado. A \omterm{def:g6:measure:area}{área} dele, passo a
passo:
\begin{enumerate}
  \item retângulo inteiro: $6 \times 4 = 24$ m$^2$;
  \item quadrado retirado: $2 \times 2 = 4$ m$^2$;
  \item \omterm{def:g6:measure:area}{área} que sobra: $24 - 4 = 20$ m$^2$.
\end{enumerate}
O \emph{\omterm{def:g6:measure:perimeter}{perímetro}} dele não é $20$ de coisa nenhuma: percorrendo o L, o
contorno continua medindo $6 + 4 + 6 + 4 = 20$ m --- os dois cortes do
canto substituem dois pedaços iguais de parede. O mesmo número por
coincidência, mas metros quadrados num caso e metros no outro!
\end{example}

\begin{remark}[Mesmo perímetro, áreas diferentes]
Duas figuras podem ter o mesmo \omterm{def:g6:measure:perimeter}{perímetro} e \omterm{def:g6:measure:area}{áreas} bem diferentes: um
quadrado $5 \times 5$ e um retângulo $9 \times 1$ têm ambos \omterm{def:g6:measure:perimeter}{perímetro}
$20$, mas \omterm{def:g6:measure:area}{áreas} $25$ e $9$. \omterm{def:g6:measure:perimeter}{Perímetro} e \omterm{def:g6:measure:area}{área} são grandezas realmente
independentes.
\end{remark}

\section{Volumes}

\begin{definition}[Volume]\label{def:g6:measure:volume}
O \emph{volume}\index{volume} de um \omterm{def:g5:solids:def}{sólido} mede o espaço que ele
preenche: quantos \omterm{def:g5:solids:def}{cubos} unitários cabem dentro dele. Unidades: cm$^3$,
m$^3$, \dots\ com $1$ m$^3 = 1\,000\,000$ cm$^3$ (cada degrau da escada
vale $1000$). Para líquidos usa-se também o \emph{litro}:
\[
1 \text{ L} = 1 \text{ dm}^3 = 1000 \text{ cm}^3,
\qquad
1 \text{ m}^3 = 1000 \text{ L}.
\]
\end{definition}

\begin{proposition}[Volume de um bloco retangular]\label{prop:g6:measure:box}
Um \omterm{def:g5:solids:def}{bloco retangular} (um \emph{prisma de base retangular}) de
comprimento $L$, largura $w$ e altura $h$ tem \omterm{def:g6:measure:volume}{volume}
\[
V = L \times w \times h .
\]
\end{proposition}

\begin{proof}
A camada de baixo contém $L \times w$ \omterm{def:g5:solids:def}{cubos} unitários (o desenho da
\cref{def:g6:measure:area}, com \omterm{def:g5:solids:def}{cubos}), e há $h$ camadas dessas.
\end{proof}

\begin{omfigure}
\begin{tikzpicture}[scale=0.6]
  \foreach \y in {0,1} {
    \foreach \x in {0,...,3} {
      \foreach \z in {0,1} {
        \begin{scope}[shift={({\x+0.45*\z},{\y+0.3*\z})}]
          \draw[fill=omDef!12] (0,0) rectangle (1,1);
          \draw[fill=omDef!20] (0,1) -- (0.45,1.3) -- (1.45,1.3) -- (1,1)
            -- cycle;
          \draw[fill=omDef!28] (1,0) -- (1.45,0.3) -- (1.45,1.3) -- (1,1)
            -- cycle;
        \end{scope}
      }
    }
  }
  \node[below, font=\small] at (2.5,-0.3)
    {$4 \times 2 \times 2 = 16$ cubos unitários};
\end{tikzpicture}

{\small Um bloco preenchido com \omterm{def:g5:solids:def}{cubos} unitários: $4 \times 2$ \omterm{def:g5:solids:def}{cubos} por
camada, duas camadas.}
\end{omfigure}

\begin{example}\label{ex:g6:measure:aquarium}
Um aquário mede $60$ cm por $30$ cm por $40$ cm (altura). \omterm{def:g6:measure:volume}{Volume}:
\[
V = 60 \times 30 \times 40 = 72\,000 \text{ cm}^3 = 72 \text{ dm}^3
= 72 \text{ L}.
\]
Passo a passo da conversão: $1000$ cm$^3$ fazem um litro, e
$72\,000 \div 1000 = 72$.
\end{example}

\section{Exercícios}

\begin{exercise}[$\star$]\label{exo:g6:measure:1}
Converta: $3.5$ m em cm; $420$ mm em cm; $0.8$ km em m;
$25\,000$ m em km.
\end{exercise}

\begin{exercise}[$\star$]\label{exo:g6:measure:2}
Calcule o \omterm{def:g6:measure:perimeter}{perímetro} de: um retângulo $8$ cm $\times$ $3.5$ cm; um
quadrado de lado $6.2$ cm; um triângulo de lados $5$ cm, $7$ cm e
$9$ cm.
\end{exercise}

\begin{exercise}[$\star$]\label{exo:g6:measure:3}
Uma pista de corrida circular tem \omterm{def:g3:shapes:shapes}{raio} $50$ m. Que comprimento tem uma
volta (valor exato com $\pi$ e depois arredondado ao metro)? Quantas
voltas fazem pelo menos $2$ km?
\end{exercise}

\begin{exercise}[$\star$]\label{exo:g6:measure:4}
Calcule a \omterm{def:g6:measure:area}{área} de: um retângulo $9$ cm $\times$ $4$ cm; um quadrado de
lado $7$ m; um \omterm{def:g6:shapes:triangles}{triângulo retângulo} de catetos $6$ cm e $10$ cm; um
disco de \omterm{def:g3:shapes:shapes}{raio} $3$ cm (valor exato e depois arredondado ao cm$^2$).
\end{exercise}

\begin{exercise}[$\star$]\label{exo:g6:measure:5}
Converta: $3$ m$^2$ em cm$^2$; $45\,000$ cm$^2$ em m$^2$; $2.5$ L em
cm$^3$; $4\,500$ L em m$^3$.
\end{exercise}

\begin{exercise}[$\star$]\label{exo:g6:measure:6}
Um campo retangular tem $120$ m de comprimento e $85$ m de largura.
Quantos metros de cerca são necessários para cercá-lo? Qual é a \omterm{def:g6:measure:area}{área}
dele?
\end{exercise}

\begin{exercise}[$\star$]\label{exo:g6:measure:7}
Calcule o \omterm{def:g6:measure:volume}{volume} de um \omterm{def:g5:solids:def}{bloco retangular}
$5$ cm $\times$ $4$ cm $\times$ $10$ cm e o de um \omterm{def:g5:solids:def}{cubo} de aresta
$3$ cm.
\end{exercise}

\begin{exercise}[$\star$]\label{exo:g6:measure:8}
Trace dois retângulos diferentes de \omterm{def:g6:measure:perimeter}{perímetro} $16$ cm e calcule as
\omterm{def:g6:measure:area}{áreas} deles. Qual dos seus retângulos tem a \omterm{def:g6:measure:area}{área} maior?
\end{exercise}

\begin{exercise}[$\star\star$]\label{exo:g6:measure:9}
Uma figura em forma de T é feita de um retângulo horizontal
$10 \times 2$ em cima de um vertical $2 \times 6$ (medidas em cm).
Calcule a \omterm{def:g6:measure:area}{área} dela e depois o \omterm{def:g6:measure:perimeter}{perímetro} (percorra o contorno com
cuidado, somando cada borda).
\end{exercise}

\begin{exercise}[$\star\star$]\label{exo:g6:measure:10}
Uma piscina é um \omterm{def:g5:solids:def}{bloco retangular} de $25$ m de comprimento, $10$ m de
largura e $2$ m de profundidade.
\begin{enumerate}
  \item Quantos metros cúbicos de água ela leva quando cheia?
  \item Quantos litros são?
  \item A piscina é enchida a $50\,000$ L por hora. Quanto tempo leva o
        enchimento?
\end{enumerate}
\end{exercise}

\begin{exercise}[$\star\star$]\label{exo:g6:measure:11}
Um jardim é um quadrado de lado $20$ m que contém um lago circular de
\omterm{def:g3:shapes:shapes}{raio} $3$ m. Que \omterm{def:g6:measure:area}{área} de grama existe ali (valor exato e depois
arredondado ao m$^2$)?
\end{exercise}

\begin{exercise}[$\star\star\star$]\label{exo:g6:measure:12}
Uma barra de chocolate mede $15$ cm $\times$ $6$ cm $\times$ $1$ cm. O
fabricante dobra \emph{as três} dimensões.
\begin{enumerate}
  \item Por quanto o \omterm{def:g6:measure:volume}{volume} fica multiplicado? Verifique calculando os
        dois \omterm{def:g6:measure:volume}{volumes}.
  \item O preço é multiplicado por $4$. A barra grande vale mais a pena
        que a pequena?
\end{enumerate}
\end{exercise}

\section{Problema: a cerca da rainha Dido}

\begin{problem}[{Problema de fim de semana --- com um comprimento fixo
de cerca, que forma cerca mais terra? O quadrado vence qualquer
retângulo, o círculo vence o quadrado, e a parede de um celeiro muda
tudo}]
\label{pb:g6:measure:1}
A lenda conta que a rainha Dido, ao desembarcar na costa da África,
recebeu ``tanta terra quanto um couro de boi pudesse cercar'' --- e ela
cortou o couro numa tira imensamente longa e fina, cercando chão
suficiente para fundar a cidade de Cartago. O problema dela agora é
seu: um agricultor tem exatamente $20$ m de cerca. O
\cref{exo:g6:measure:8} mostrou que dois currais de mesmo \omterm{def:g6:measure:perimeter}{perímetro}
podem ter \omterm{def:g6:measure:area}{áreas} diferentes; este problema encontra o \emph{melhor}
curral --- e descobre que a resposta muda completamente quando a parede
de um celeiro empresta um lado de graça.

\medskip
\noindent\textbf{Parte I --- Vinte metros de cerca.}
\begin{enumerate}
  \item O curral tem de ser um retângulo usando todos os $20$ m de
        cerca. Confira que um curral $1 \times 9$ e um $2 \times 8$
        servem, e calcule as \omterm{def:g6:measure:area}{áreas} deles.
  \item Explique por que o comprimento e a largura de todo curral que
        serve somam $10$ m. Depois monte a tabela completa dos currais
        de lados inteiros ($1 \times 9$ até $5 \times 5$) com as \omterm{def:g6:measure:area}{áreas}
        deles. Qual é o melhor?
  \item Vale a pena testar lados decimais? Calcule as \omterm{def:g6:measure:area}{áreas} do curral
        $4.5 \times 5.5$ e do curral $4.9 \times 5.1$, e compare com o
        quadrado $5 \times 5$. O que você conjectura?
  \item A questão 2 transformou o problema da cerca numa questão
        puramente numérica: \emph{entre todos os pares de números que
        somam $10$, que par tem o maior \omterm{def:g2:mult:def}{produto}?} Responda a partir da
        sua tabela e enuncie a regra geral que você observa.
  \item Eis por que afastar-se do quadrado sempre faz perder, com
        tesoura em vez de álgebra: parta do curral $5 \times 5$ e
        transforme-o no curral $6 \times 4$ retirando uma faixa e
        colando outra de volta. Que faixa é retirada, qual é
        acrescentada, e por que a troca perde exatamente um metro
        quadrado? Explique por que um passo a mais (até $7 \times 3$)
        perde ainda mais.
\end{enumerate}

\medskip
\noindent\textbf{Parte II --- O celeiro e o círculo.}
\begin{enumerate}[resume]
  \item Agora o curral é construído encostado na parede comprida de um
        celeiro: a parede substitui um dos comprimentos do curral, e os
        $20$ m de cerca cobrem apenas os outros três lados. Monte a
        tabela dos currais de lados inteiros (largura $w$, comprimento
        $L$ ao longo da parede, $w + L + w = 20$) e as \omterm{def:g6:measure:area}{áreas} deles. Que
        curral ganha agora --- e ele é um quadrado?
  \item Explique o vencedor com um espelho (\cref{ch:g6:symmetry}):
        reflita o curral em relação à parede do celeiro e considere o
        curral duplicado. Quanta cerca o curral duplicado usa, qual
        curral duplicado é o melhor segundo a Parte I, e o que isso faz
        do curral original?
  \item Dido não construiu um retângulo. Dobre os $20$ m de cerca num
        círculo: usando $\pi \approx 3.14$, calcule o \omterm{def:g3:shapes:shapes}{raio} dele (ao
        centímetro) e depois a \omterm{def:g6:measure:area}{área} (ao m$^2$)
        (\cref{prop:g6:measure:circle},
        \cref{prop:g6:measure:formulas}). Compare com o quadrado
        $5 \times 5$.
  \item O problema inverso: quer-se um curral de \omterm{def:g6:measure:area}{área} exatamente
        $36$ m$^2$, com o mínimo possível de cerca. Compare os
        retângulos $1 \times 36$, $2 \times 18$, $3 \times 12$,
        $4 \times 9$ e $6 \times 6$: \omterm{def:g6:measure:perimeter}{perímetros}? Que forma é a mais
        barata, e como esta questão é a imagem espelhada da Parte I?
  \item Em uma frase: por que latas, canos e caixas d'água são tantas
        vezes \emph{redondos}?
\end{enumerate}

\medskip
\noindent\textbf{Parte III --- \omterm{def:g6:measure:perimeter}{Perímetro} e \omterm{def:g6:measure:area}{área} são estranhos um ao
outro.}
\begin{enumerate}[resume]
  \item Encontre um retângulo cujo \omterm{def:g6:measure:perimeter}{perímetro} passe de $100$ m mas cuja
        \omterm{def:g6:measure:area}{área} seja menor que $1$ m$^2$. (Bem comprido e bem fino.
        Decimais permitidos.)
  \item Verdadeiro ou falso: ``de duas figuras, a de maior \omterm{def:g6:measure:perimeter}{perímetro}
        tem maior \omterm{def:g6:measure:area}{área}''. Dê um contraexemplo tirado deste problema.
  \item Pegue o retângulo $6 \times 4$ e recorte um entalhe quadrado
        $1 \times 1$ no meio de um dos lados compridos (o entalhe se
        abre para fora, como um dente que faltasse). Calcule a \omterm{def:g6:measure:area}{área} e o
        \omterm{def:g6:measure:perimeter}{perímetro} da figura entalhada e compare os dois com os do
        original. O que aconteceu com cada um?
  \item Os cartógrafos conhecem um fato estranho: quanto mais detalhado
        o mapa, \emph{mais comprido} um litoral fica --- cada ampliação
        revela novos entalhezinhos, e a questão 13 mostra o que cada
        entalhe faz. Explique em uma ou duas frases por que ``o
        comprimento do litoral da Bretanha'' é um número escorregadio,
        enquanto ``a \omterm{def:g6:measure:area}{área} da Bretanha'' não é.
  \item A prova do agricultor. Com $36$ m de cerca e a parede do
        celeiro disponível, encontre o melhor curral retangular (use o
        espelho da questão 7), dê a \omterm{def:g6:measure:area}{área} dele e compare com o melhor
        curral construído com $36$ m em campo aberto. Quanto a parede
        do celeiro rende ao agricultor?
\end{enumerate}
\end{problem}
